\subsection{\texorpdfstring{\(\mathbf{R}^n\) 的离散子群}{R 的 n 次幂的离散子群}}

我们在第 V 章（§ 1 第 1 小节命题 1）中已经看到，R 的闭子群中除 R 自身外只有离散子群，它们都由单个元素生成。我们将从考虑 \(\mathbf{R}^n\) 的离散子群开始。

首先，由 \(\mathbf{R}^n\) 的典范基（第 VI 章 § 1 第 3 小节）的 \(p\) 个向量（ \(p \leqslant n\) ）所生成的 \(\mathbf{R}^n\) 的子群是一个离散群，同构于乘积 \(\mathbf{Z}^p\) ——即 \(p\) 个都等于 \(\mathbf{Z}\) 的群的乘积。更一般地，考虑子群 \(\mathbf{G}\) ，它由构成自由系的 \(p\) 个点 \(a_i\) （ \(i \leqslant i \leqslant p\) ）生成。存在 \(\mathbf{R}^n\) 到其自身上的一个双射线性映射，它把 \(a_i\) 映为 \(e_i\) （ \(i \leqslant i \leqslant p\) ）；由于这样的映射是拓扑群 \(\mathbf{R}^n\) 的自同构， \(\mathbf{G}\) 作为拓扑群同构于由 \(e_i\) （ \(i \leqslant i \leqslant p\) ）生成的子群，因此是一个秩为 \(p\) 、同构于 \(\mathbf{Z}^p\) 的离散子群。

群 \(Z^{p}\) 的结构、从而 G 的结构，已在代数学中研究过。我们回顾这一研究的主要结果。G 关于环 Z 的基是由 p 个点组成的系

\[
\boldsymbol {b} _ {i} = \sum_ {j = 1} ^ {p} r _ {i j} \boldsymbol {a} _ {j},
\]

其中 \(r_{ij}\) 是使得行列式 \(\det(r_{ij})\) 等于 +1 或 -1 的整数。G 的每个子群 H 都是离散的且秩 \(q \leqslant p;\) 此外，若 \(H\) 是给定的秩为 \(q\) 的子群，则存在由 \(p\) 个点 \(b_i\) （ \(i \leqslant i \leqslant p\) ）组成并生成 \(G\) 的自由系，以及由 \(q\) 个点 \(c_i\) （ \(i \leqslant i \leqslant q\) ）组成并生成 \(H\) 的系，使得有 \(c_i = e_i b_i\) （对 \(i \leqslant i \leqslant q\) ），其中 \(e_i\) 是整数（ \(H\) 关于 \(G\) 的不变因子），满足

\[
e _ {i + 1} \equiv o (\mathrm{mod} e _ {i}) \quad \text { for } \quad i \leqslant i \leqslant q - i.
\]

商群 G/H 是一个离散群，同构于 \(Z^{p-q} \times F\) ，其中 F 是一个有限交换群，是 q 个循环子群的直积，这些循环子群的阶分别为 \(e_{1}, \ldots, e_{q}\) 。

现在我们将证明，刚才所考虑的 \(\mathbf{R}^n\) 的离散子群是仅有的离散子群。

命题 1. 设 G 是 \(\mathbf{R}^n\) 的一个秩为 \(p\) 的离散子群， \((\boldsymbol{a}_i)_{1 \leqslant i \leqslant p}\) 是 G 的 \(p\) 个点组成的自由系，P 是以 o 为中心、以 \(\boldsymbol{a}_i\) 为基向量的闭超平行体（第 VI 章 § 1 第 3 小节）。则集合 G ∩ P 是有限的且生成 G，且 G 的每个点都是 \(\boldsymbol{a}_i\) 以有理数为系数的线性组合。

G ∩ P 是紧的且离散的，因而是有限的。设 x 是 G 的任意点；它等于线性组合 \(\sum_{i=1}^{p} t_{i}a_{i}\) （诸 \(a_{i}\) 以实数为系数）。对每个整数 m \textgreater{} 0 ，考虑点

\[
z _ {m} = m x - \sum_ {i = 1} ^ {p} [ m t _ {i} ] a _ {i} = \sum_ {i = 1} ^ {p} (m t _ {i} - [ m t _ {i} ]) a _ {i} (*);
\]

它属于 G，且由于 \(o \leqslant mt_{i} - [mt_{i}] < 1\) ，它位于 P 中。因此，首先有 \(x = z_{1} + \sum_{i=1}^{p} [t_{i}] a_{i}\) ，故 G 由 \(G \cap P\) 生成；其次，由于 \(G \cap P\) 是有限的，存在两个不同的整数 h、k，使得 \(z_{h} = z_{k}\) ，从而 \((h - k)t_{i} = [ht_{i}] - [kt_{i}]\) （对 \((1 \leqslant i \leqslant p)\) ），因此 \(t_{i}\) 是有理数。

推论. 设 \((\pmb{a}_i)_{1\leqslant i\leqslant p}\) 是 \(p\) 个点（ \(\mathbf{R}^n\) 的）组成的自由系，并设

\[
\boldsymbol {b} = \sum_ {i = 1} ^ {p} t _ {i} \boldsymbol {a} _ {i}
\]

是 \(a_{i}\) 以实数为系数的一个线性组合。则 \(R^{n}\) 中由 \(p + i\) 个点 \(a_{1}, a_{2}, \ldots, a_{p}\) 与 b 所生成的子群 G 是离散的，当且仅当诸数 \(t_{i}\) 是有理数。

命题 1 表明这条件是必要的。它也是充分的，因为若这条件满足，我们就可以写 \(t_i = m_i / d\) ，其中 \(d\) 与诸 \(m_i\) 是整数（ \(i \leqslant i \leqslant p\) ）；于是 \(b\) 是 \(p\) 个点 ( \(\mathrm{I} / d\) ) \(a_i\) 以整数为系数的线性组合；因此 \(G\) 是由这 \(p\) 个点生成的离散群的一个子群，从而本身是离散的。

命题 1 的结果可以表述如下：若 q 个点 \(x_{i}\) （ \(i \leqslant i \leqslant q\) ，属于 \(R^{n}\) 的一个离散子群 G）构成一个关于 R 相关的系，则它们构成一个关于 Q 相关的系。由此立即推出， \(R^{n}\) 的离散子群的有理秩等于它的秩。

把命题 1 的推论应用到 \(a_{i}\) 为 \(n\) 个向量 \(\pmb{e}_{i}\) （ \(\mathbf{R}^{n}\) 的典范基的）的情形，就得到下面的命题：

命题 2（克罗内克）. 设 \(\theta_{1}, \theta_{2}, \ldots, \theta_{n}\) 是 \(n\) 个实数。为了对每个 \(\varepsilon > 0\) ，存在一个整数 \(q\) 与 \(n\) 个整数

\[
p _ {i} \quad (\mathrm{I} \leqslant i \leqslant n)
\]

使得

\[
\left| q \theta_ {i} - p _ {i} \right| \leqslant \varepsilon \quad (\mathrm{I} \leqslant i \leqslant n),
\]

其中这些不等式中至少有一个的左端不为零，其充分必要条件是诸 \(\theta_{i}\) 中至少有一个是无理数。

定理 1. \(\mathbf{R}^n\) 的每个秩等于 \(p\) 的离散子群 G 都由 \(p\) 个点组成的自由系生成。

根据前面回顾过的同构于 \(Z^{p}\) 的群的性质，只需证明 G 是由 p 个点组成的自由系生成的离散群的一个子群。现在，由于 G 的秩为 p，存在 G 的由 p 个点 \(a_{i}\) （ \(1 \leqslant i \leqslant p\) ）组成的自由系，使得每个 \(x \in G\)

都等于线性组合 \(\sum_{i=1}^{p} t_i a_i\) （诸 \(a_i\) 以实数为系数）；

由于 G 是离散的，命题 1 表明诸 \(t_i\) 是有理数。此外，命题 1 表明 G 由有限多个点生成；由于这些点是 \(\boldsymbol{a}_i\) 以有理数为系数的线性组合，故存在整数 \(d\) ，使得它们是 \(p\) 个点 \((\mathrm{I}/d)\boldsymbol{a}_i = \boldsymbol{a}_i'\) 以整数为系数的线性组合。由此推出，G 是由诸 \(\boldsymbol{a}_i'\) 生成的群的一个子群。

定理 1 也可以不用不变因子理论来证明（参见习题 1）。

\(R^{n}\) 中秩为 n 的离散子群也称为 \(R^{n}\) 中的格。

